%----------------------------------------------------------------------------------------
% METADATA BLOCK (YAML)  — see inbox/README.md for every field
%----------------------------------------------------------------------------------------
% uid: ""
% school: "lycee"
% level: "4eme"
% section: "Mathématiques"
% chapter: "Coniques"
% title: "Un triangle équilatéral inscrit dans l'hyperbole"
% difficulty: "Moyen"
% tags: ["hyperbole", "triangle-equilateral", "orthocentre", "branches"]
% provenance: "adapted"
% batch: "2026-11-20-defis-geometrie-A"
% defi: true
%----------------------------------------------------------------------------------------

% ---STATEMENT---
Le plan est rapporté à un repère orthonormé. On note $\Gamma$ l'hyperbole d'équation $xy=1$ ; elle a deux branches, celle des abscisses positives et celle des abscisses négatives.
Soient $a$, $b$, $c$ trois réels non nuls deux à deux distincts et $A\left(a,\dfrac1a\right)$, $B\left(b,\dfrac1b\right)$, $C\left(c,\dfrac1c\right)$ trois points de $\Gamma$.

\begin{enumerate}
  \item Montrer que l'orthocentre $H$ du triangle $ABC$ est le point de coordonnées $\left(-\dfrac1{abc},-abc\right)$ et qu'il appartient à $\Gamma$.
  \item On suppose que $ABC$ est équilatéral, de centre de gravité $G$. Montrer que \textbf{le centre $G$ du triangle appartient à $\Gamma$}.
  \item
  \begin{enumerate}
    \item Montrer que \textbf{les trois sommets ne peuvent pas être sur la même branche de $\Gamma$}.
    \item Vérifier que $A(1,1)$, $B(-2+\sqrt3,\,-2-\sqrt3)$ et $C(-2-\sqrt3,\,-2+\sqrt3)$ sont les sommets d'un triangle équilatéral inscrit dans $\Gamma$. Calculer son centre et retrouver les résultats des questions 2 et 3a.
  \end{enumerate}
\end{enumerate}

% ---SOLUTION---
\begin{enumerate}[itemsep=14pt]
  \item\leavevmode
  \begin{prop}[Perpendicularité]
  Deux droites de coefficients directeurs $m$ et $m'$ sont perpendiculaires si et seulement si \ib{$mm'=-1$}.
  \end{prop}
  Coefficient directeur de $(BC)$ : $\dfrac{\dfrac1c-\dfrac1b}{c-b}=-\dfrac1{bc}$.

  La hauteur issue de $A$ a pour coefficient directeur $bc$ et passe par $A$ : $y-\dfrac1a=bc\,(x-a)$.

  Avec $H\left(-\dfrac1{abc},-abc\right)$ :
  \[y_H-\dfrac1a=-abc-\dfrac1a\qquad\text{et}\qquad bc(x_H-a)=-\dfrac1a-abc.\]
  Les deux nombres sont égaux : $H$ est sur la hauteur issue de $A$. Les coordonnées de $H$ sont symétriques en $a,b,c$ : il est aussi sur les deux autres hauteurs, donc \ib{$H$ est l'orthocentre}.

  Enfin $x_Hy_H=1$ : $H\in\Gamma$.

  \item\leavevmode
  \begin{prop}[Triangle équilatéral]
  Dans un triangle équilatéral, chaque hauteur est aussi médiane (et médiatrice). Orthocentre et centre de gravité sont donc confondus.
  \end{prop}
  Ici $G=H$. D'après la question 1, \ib{$G=H\in\Gamma$}.

  \item
  \begin{enumerate}[itemsep=8pt]
    \item $G$ a pour coordonnées $\left(\dfrac{a+b+c}3,\ \dfrac{\frac1a+\frac1b+\frac1c}3\right)$ et $G=H$, donc
    \[\dfrac{a+b+c}{3}=-\dfrac{1}{abc}.\]
    \emph{Cas 1 : $a,b,c>0$.} Le membre de gauche est $>0$ ; le membre de droite est $<0$. Impossible.

    \emph{Cas 2 : $a,b,c<0$.} Le membre de gauche est $<0$ ; comme $abc<0$, le membre de droite $-\dfrac1{abc}$ est $>0$. Impossible.

    Donc $a$, $b$, $c$ n'ont pas tous le même signe : les trois sommets ne sont pas sur la même branche.
    \item \emph{Sur $\Gamma$ :} $1\times1=1$ et $(-2+\sqrt3)(-2-\sqrt3)=4-3=1$ ; $B$ et $C$ sont donc sur $\Gamma$ (même produit $x\,y$ avec $x,y$ échangés pour $C$).

    \emph{Côtés :}
    \[AB^2=(3-\sqrt3)^2+(3+\sqrt3)^2=12-6\sqrt3+12+6\sqrt3=24,\]
    \[BC^2=(2\sqrt3)^2+(2\sqrt3)^2=24,\qquad CA^2=(3+\sqrt3)^2+(3-\sqrt3)^2=24.\]
    Les trois côtés valent $\sqrt{24}$ : $ABC$ est équilatéral.

    \emph{Centre :} $G\left(\dfrac{1-2-2}{3},\dfrac{1-2-2}{3}\right)=(-1,-1)$ et $(-1)(-1)=1$ : $G\in\Gamma$ (question 2).

    \emph{Branches :} $A$ est sur la branche $x>0$, $B$ et $C$ sur la branche $x<0$ (comme $G$) : les sommets ne sont pas tous sur la même branche (question 3a).
  \end{enumerate}
\end{enumerate}

\textbf{Idée :} tout vient de $G=H$ pour un triangle équilatéral, et du fait que l'orthocentre d'un triangle inscrit dans $xy=1$ est connu en fonction de $a,b,c$.

\textbf{Piège :} $G$ est le centre de gravité (moyenne des sommets), pas le centre du cercle circonscrit ; pour un triangle équilatéral c'est le même point, mais il faut le dire avant de poser $G=H$.
